% main.tex -- Do categorial diagnostics measure anything?
% House style: CLAUDE.md. Outcome-neutrality rule: PLAN.md section 5.
% ⟨RESULT⟩ = number pending from results/; ⟨VERIFY⟩ = claim pending
% source check (see TO_VERIFY.md). Neither may be resolved by the
% executor from memory.

\documentclass[12pt]{article}

% House style preamble (snapshot of .house-style/preamble.tex; XeLaTeX).
% Loads fonts, biblatex-APA, csquotes, langsci-gb4e, semantic macros, and
% \addbibresource{references.bib} (the verified-only project bib).
\input{.house-style/preamble.tex}

\title{Do categorial diagnostics measure anything?\\
\large A psychometric study of English lexical categories}
\author{Brett Reynolds\\Humber Polytechnic \& University of Toronto}
\date{}

\begin{document}
\maketitle

\begin{abstract}
\noindent ⟨RESULT: 150-word abstract after Phase 4; commits to reporting
whichever measurement model wins per the pre-registered comparison⟩
\end{abstract}

\section{Introduction}\label{sec:intro}

Syntacticians settle lexical-category membership with diagnostics. Does the word
take \textit{very}? Does it license an NP complement? Does it occur
attributively? For a century, these tests have functioned as the field's
measurement instruments. Like any instruments, they presuppose
answers to questions nobody has asked of them. Do they cohere? What does
each one actually indicate? When two diagnostics disagree about the same
word, which one is broken?

Psychology confronted these questions about its own instruments decades
ago and built a discipline (psychometrics) around them. Linguistics
didn't. The diagnostic toolkit passed from textbook to textbook with its
validity assumed rather than estimated. This paper estimates it.

The empirical object is a lexeme\,$\times$\,diagnostic acceptability
matrix: rows are words, columns are standard categorial diagnostics
instantiated as test sentences, cells are acceptability ratings. Given
such a matrix, the question \enquote{do these diagnostics measure
anything?} becomes statistically precise. If the diagnostics for, say,
adjectivehood jointly reflect one underlying property, the matrix should
be well described by a model with few latent dimensions: the diagnostics
behave as items on a test, and the category behaves as the construct the
test measures. If instead the diagnostics form a web of mutual
constraints with no common cause (each predicting some neighbours,
none reflecting a shared essence), the matrix should be better described
by a sparse network model. These are competing measurement models, and
competing measurement models can be fitted and compared.

The contrast matters because it operationalizes a standing theoretical
dispute. The traditional view treats categories as discrete, with
diagnostics detecting membership in them. \textcite{croft2001} argues the
diagnostics are opportunistically selected and define no coherent
categories at all. \textcite{aarts2007} keeps category boundaries sharp
but grades prototypicality within them. His subsective gradience makes
\textit{utter} as much an adjective as \textit{thin}, just a less
prototypical one. A fourth position, the one this paper's framework comes
from, holds that categories are clusters of mutually reinforcing
properties, maintained by mechanisms rather than defined by essences.
Each position predicts a different matrix structure
(\S\ref{sec:predictions}), so the data can adjudicate. I commit in
advance: the analysis plan, model
comparison, and decision thresholds were fixed before the data were
analyzed, and this introduction is written to survive any outcome,
because which model wins \emph{is} the finding.

A category label, on any of these views, earns its keep by licensing
inference. Call a categorization \emph{projectible}, for a given purpose,
to the extent that knowing an item's category supports reliable
predictions about its unexamined behaviour, a term I adapt from
\textcite{goodman1955}. Projectibility gives the validity question teeth
at the level of the individual diagnostic: a test that predicts nothing
beyond itself is decoration, whatever its pedigree. I estimate a
projectibility score for every diagnostic in the battery
(\S\ref{sec:methods}): the improvement in out-of-sample prediction of a
word's remaining profile that comes from observing that one test.

The paper makes three contributions. First, a validated method: I show
on data with known structure that the model comparison discriminates
reflective from network organization, so a verdict on real data is
licensed rather than assumed. Second, a reanalysis of the largest
existing acceptability grid, the MegaAcceptability dataset of 1{,}000
clause-embedding verbs in 50 frames \parencite{whiterawlins2016},
asking the measurement question its creators didn't: ⟨RESULT⟩. Third, a
new dense grid over the adjective--preposition--determinative region,
including the boundary items the descriptive literature has long
flagged (\textit{worth}, \textit{near}, \textit{fun}, \textit{galore}),
with quantified rather than anecdotal misfit: ⟨RESULT⟩.

\section{Background}\label{sec:background}

\subsection{The diagnostic tradition and its discontents}

Distributional diagnostics are as old as structuralism, but their
troubles are best seen at category boundaries. \textit{Worth} licenses
an NP complement like a preposition (\textit{worth a fortune}), resists
\textit{very} (\textit{*very worth the trip}) while accepting
\textit{well} (\textit{well worth the trip}), inverting the pattern of
central adjectives, and occurs as complement
of \textit{become} like an adjective, with a gapped gerund-participial
complement (\textit{worth reading \underline{\hspace{1em}}}) that
clusters with adjectival \textit{tough}-constructions.
\textcite{maling1983} read the balance of evidence as completed
reanalysis to preposition; the \textit{Cambridge Grammar} weighs the
same conflicts and keeps \textit{worth} in the adjective category,
judging the case strong, and conducts the argument in its
prepositions chapter, which rather concedes the point about where the
word lives \parencite[ch.~7, \S2.2]{cgel2002}. The disagreement isn't
peculiar to \textit{worth}. The field's two most detailed treatments
invert each other on both flagship items: \textcite{maling1983} makes
\textit{worth} a preposition and \textit{near} an adjective
(\enquote{the only surviving relic of the class of transitive
adjectives}), while the \textit{Cambridge Grammar} makes \textit{worth}
an adjective and \textit{near} an exceptional preposition, and
\textcite{aarts2007} splits \textit{near} by syntactic configuration.
Same data, opposite assignments: the diagnostics underdetermine the
verdict. The point here isn't
who's right. It's that both verdicts discard stable, replicable facts
as noise, because the adjudication game has no role for them.

\textcite{matthews2014} documents a second discontent: the criteria
themselves were quietly demoted. The discipline retreated from
necessary-and-sufficient definitions of the adjective category to a
\enquote{cluster of syntactic properties} characterizing central
members, with positional evidence at the boundaries conceded to be
\enquote{less than decisive}, which is to say, the field already
treats its categories as clusters and its diagnostics as fallible
indicators; it just hasn't measured either. And \textcite{croft2001} presses the radical version of the
complaint: since no two diagnostics pick out exactly the same category,
the analyst's choice of which to trust is opportunistic, and the
categories so established are artefacts of method. Croft's argument is
usually met with rebuttal by example. It's better met with measurement:
opportunism is only a vice if the diagnostics genuinely fail to cohere,
and coherence is an estimable quantity.

\subsection{Categories as instruments}

Treating diagnostics as test items makes the dispute between these
positions empirical. In a \emph{reflective} measurement model, a latent
variable (the construct) causes the item responses; items correlate
\emph{because} they share a cause, and the model implies a low-rank
structure on the response matrix. In a \emph{network} model, no common
cause exists: items constrain one another directly, pairwise, and the
implied structure is a sparse graph. The contrast has restructured
measurement debates in psychology~\parencite{borsboomcramer2013},
and it maps onto the theoretical landscape with almost embarrassing
neatness: essence-detection is reflective; the maintained-cluster view
is a structured network; Croftian fragmentation is a degenerate graph
with no stable communities; Aartsian gradience is reflective structure
read through graded scores.

One technical caution governs everything that follows: Ising network
models and item-response models can be statistically
equivalent (descriptions of the same data distribution), so
goodness of fit alone licenses no inference about which structure
generated the data \parencite{marsman2018}, and low-rank approximation
adds a further idealization downstream of that equivalence. This is
why the comparison here is never raw fit: I compare structural
commitments under matched complexity (held-out prediction, parsimony,
and structure diagnostics together; \S\ref{sec:methods}) and I
verify on planted structure that the joint battery discriminates the
regimes at the sample sizes.

\subsection{Existing judgment data}

⟨Drafted in Phase 2 after the citing-papers sweep: MegaAcceptability and
the bleaching method \parencite{whiterawlins2016, whiterawlins2020};
Sprouse--Almeida replication datasets as the hard-region validation set;
CoLA; what none of them asked.⟩

\section{Predictions}\label{sec:predictions}

Each position implies a different structure for the
lexeme\,$\times$\,diagnostic matrix, so the data can adjudicate between them.
The mapping in Table~\ref{tab:predictions}, and the decision rule that follows
it, were registered before any rating existed. Which structure the data show is
the finding, and this section is written to read the same whichever way it
falls.

\begin{table}[t]
\centering\small
\caption{What each position predicts for the lexeme\,$\times$\,diagnostic
matrix, in plain terms. \emph{Predicts better}: which model more accurately fills
in ratings it was not shown. \emph{Dominant dimension}: whether one underlying
dimension towers over the rest (the fingerprint of a single common cause), the
variation spreads across several with no leader (mutual constraint,
\emph{not dominant}), or there is no structure at all (\emph{flat}).
\emph{Clusters}: whether words grouped by their diagnostic profiles recover the
traditional categories. \emph{Boundary items}: how the pre-named hard cases
(\textit{worth}, \textit{near}, and the rest) behave under the winning model.}
\label{tab:predictions}
\begin{tabular}{@{}lllll@{}}
\toprule
Position & Predicts better & Dominant dimension & Clusters & Boundary items \\
\midrule
Discrete essence & reflective & dominant & match categories & fit normally \\
Croftian fragmentation & neither & flat & none stable & scattered \\
Aartsian gradience & reflective & dominant & match categories & fit, score low \\
Maintained cluster & network & not dominant & categories + overlap & misfit \\
\bottomrule
\end{tabular}
\end{table}

Two positions, the discrete-essence tradition and \textcite{aarts2007}, both
predict a reflective matrix, and they part only on the boundary items. Under
discrete essence a word is in a category or out of it, so \textit{worth} and
\textit{near} should fit the reflective model as well as any central member.
Under subsective gradience they fit it too, but at low prototypicality, a graded
position inside a sharply bounded category rather than extra structure. The
maintained-cluster view predicts instead that they misfit the winning model, and
\textcite{croft2001} predicts no stable model for them to misfit at all.

The verdict turns on prediction. I compare the two models by how well each fills
in ratings it was not shown: a random tenth of the cells is hidden, each model is
fitted to what remains, and its predictions are checked against the hidden values
(cross-validation; \S\ref{sec:methods} gives the procedure). A model that has
found real structure predicts the hidden cells well; one that has only traced
noise does not. Hide the cell for \textit{worth} with \textit{very}, say: the
reflective model predicts it from where \textit{worth} sits on a shared adjective
dimension, the network model from \textit{worth}'s direct links to the other
tests, and the closer prediction wins, totalled over many hidden cells.

The two scores rarely tie, so I ask whether the gap between them is real or
chance. Each carries a margin of error, and I treat the models as genuinely
separated only when their gap is at least twice that margin, the familiar
two-standard-error bar; a narrower gap is a tie.

Prediction says which model fits; the shape of the correlations says why. If a
single common cause drives the diagnostics, their shared variation piles onto one
large dimension, as a strong first factor does in factor analysis. If the
diagnostics instead constrain each other piecemeal, that variation splits across
several comparable dimensions with no leader. I read this off the largest
dimension relative to the next (the ratio of the first two eigenvalues of the
diagnostics' association matrix): above three points to a common cause, below
three to a web of local constraints.

Two safeguards keep the verdict honest. Acceptability is a matter of degree, but
the models read a flat accept-or-reject, so each rating is first cut into a yes or
no; because that cut-off is a choice, the verdict must hold under several of them,
and one that reverses when the cut-off moves is reported as unstable, not as a
finding. And a battery is a measuring instrument only if its tests inform one
another: at least ten diagnostics must measurably sharpen prediction of a word's
remaining profile, their projectibility (\S\ref{sec:intro}), or the toolkit
measures nothing.

The boundary prediction attaches to named words. A word \emph{misfits} when the
winning model cannot reproduce its profile: its pattern of acceptances and
rejections pulls in directions the model does not expect. The pre-named set,
\textit{worth}, \textit{near}, \textit{fun}, \textit{galore}, \textit{like},
\textit{due}, \textit{ablaze}, \textit{asleep}, \textit{more}, \textit{less},
\textit{enough}, \textit{such}, and \textit{own}, is predicted to land among the
worst-fitting tenth of words (the top misfit decile); if it doesn't, the informal
boundary-item literature is miscalibrated, itself a result. Claims about any
single word are held to a floor of forty diagnostics, below which one word's
misfit is too noisy to read.

One sharper prediction is registered in advance from independent syntactic
work: the profiles of \textit{more} and \textit{less} should cluster with the
determinatives, not with the adjectives or the adverbs.
% ⟨pm: P1's more/less cite is a SEPARATE work from reynolds2021 (= "Quantifying
% determinatives", Cadernos); Brett to name the more/less paper's venue/status
% before this cite is added. See DECISIONS 2026-06-11.⟩

The commitment runs the other way too. If a clean reflective model fits, with one
dimension per category and unremarkable misfit at the boundary, the
maintained-cluster framing is wrong for this region and the discussion will say
so. Three of the four rows in Table~\ref{tab:predictions} would each falsify the
position this paper's framework comes from.

\section{Study 1: MegaAcceptability reanalysis}\label{sec:study1}
⟨RESULT⟩

\section{Study 2: the category-boundary grid}\label{sec:methods}
⟨Battery (diagnostics.yaml); bleached instantiation; LLM judge with
designed gold sample and bias-corrected estimation ⟨VERIFY: PPI/DSL
citations⟩; judge validation against published judgments.⟩

\section{Results}\label{sec:results}
⟨RESULT⟩

\section{Discussion}\label{sec:discussion}
⟨Phase 4. Must include the concession paragraph required by PLAN.md
\S5 if the clean reflective model wins.⟩

\printbibliography
\end{document}
